%% (H1) Strojové učení
\subsection{Učení s učitelem}
\begin{itemize}
\item klasifikace, regrese (základní typy úloh)
\begin{itemize}
\item úloha klasifikace spočívá v přiřazení vhodné třídy (nebo sady tříd) konkrétní položce
\begin{itemize}
\item typická je binární klasifikace
\end{itemize}
\item při regresi přiřazujeme položce $x$ vhodnou funkční hodnotu $y$
\end{itemize}
\item standardní evaluační metriky
\begin{itemize}
\item regrese
\begin{itemize}
\item mean square error, $\mathrm{MSE}=\frac1N\sum_{i=1}^N (y_i-t_i)^2$
\begin{itemize}
\item $y_i$ \dots{} predikce funkční hodnoty pro $x_i$
\item $t_i$ \dots{} reálná funkční hodnota pro $x_i$ (v trénovacích/testovacích datech)
\end{itemize}
\end{itemize}
\item klasifikace
\begin{itemize}
\item negative log likelihood, $\mathrm{NLL}=-\sum_{i=1}^N \log p_i(t_i)$
\begin{itemize}
\item $t_i$ \dots{} reálná třída $x_i$
\item $p_i$ \dots{} pravděpodobnostní rozdělení, které model vrací pro $x_i$ (pravděpodobnosti, že $x_i$ náleží jednotlivým třídám)
\item místo $p_i(t_i)$ lze také zapsat $p(t_i\mid x_i)$
\end{itemize}
\item accuracy (přesnost)
\begin{itemize}
\item počet správně klasifikovaných $\div$ celkový počet
\item $\frac{tp+tn}{tp+tn+fp+fn}$ (u binární klasifikace)
\item nevhodná metrika, pokud jsou třídy nevyvážené (např. testujeme vzácné onemocnění)
\end{itemize}
\item confusion matrix
\begin{itemize}
\item řádky jsou označeny pravdivými hodnotami, sloupce predikovanými hodnotami (nebo naopak - je to potřeba uvést)
\item v každé buňce je počet testovacích dat s konkrétní kombinací targetu a predikce
\item na diagonále jsou tedy počty správně klasifikovaných testovacích dat
\item při binární klasifikaci buňky vlastně obsahují počty TP, FN, FP, TN
\item \href{https://www.evidentlyai.com/classification-metrics/confusion-matrix#how-to-read-a-confusion-matrix}{[obr.: confusion matrix]}
\end{itemize}
\end{itemize}
\item binární klasifikace (navíc)
\begin{itemize}
\item základní dělení výsledků klasifikace
\begin{itemize}
\item predikce je pozitivní, realita (target) rovněž $\to$ true positive (TP)
\item predikce je pozitivní, realita je negativní $\to$ false positive (FP)
\item predikce je negativní, realita je pozitivní $\to$ false negative (FN)
\item predikce je negativní, realita také $\to$ true negative (TN)
\end{itemize}
\item precision, $P=\frac{tp}{tp+fp}$
\item recall, $R=\frac{tp}{tp+fn}$
\item F-measure
\begin{itemize}
\item vážený harmonický průměr precision a recallu
\item $F_\beta=\frac{(\beta^2+1)PR}{\beta^2P+R}=\frac{tp+\beta^2\cdot tp}{tp+fp+\beta^2(tp+fn)}$
\item $F_1=\frac{2PR}{P+R}=\frac{tp+tp}{tp+fp+tp+fn}$
\item je užitečné mít jednu hodnotu místo dvou ($P,R$)
\end{itemize}
\item ROC (receiver operating characteristic) křivka
\begin{itemize}
\item zachycuje, jak se mění vlastnosti binárního klasifikátoru, když posouváme práh (threshold, tzn. hodnotu, od níž výsledek prohlásíme za pozitivní - typicky uvažujeme 0.5)
\item na ose $x$ je false positive rate, $\mathrm{FPR}=\frac{fp}{fp+tn}$ (v čitateli jsou všechny negativní instance)
\item na ose $y$ je true positive rate (recall), $\mathrm{TPR}=\frac{tp}{tp+fn}$ (v čitateli jsou všechny pozitivní instance)
\item každý bod na křivce odpovídá jinému thresholdu
\item pro náhodný \quotedblbase{}klasifikátor\textquotedblleft{} bude ROC křivka diagonála
\item pro dokonalý klasifikátor bude ROC křivka jeden bod vlevo nahoře
\begin{itemize}
\item FPR = 0, TPR = 1
\end{itemize}
\end{itemize}
\item AUC (area under curve)
\begin{itemize}
\item plocha pod ROC křivkou
\item pro náhodný \quotedblbase{}klasifikátor\textquotedblleft{} je AUC = 0.5, pro dokonalý klasifikátor to bude AUC = 1
\end{itemize}
\item hodnoty metrik typu true positive rate, false negative rate, \dots{} se vždycky určují tak, že v čitateli jsou všechny instance odpovídající názvu a ve jmenovateli jsou všechny instance, co sdílejí \quotedblbase{}realitu\textquotedblleft{} s čitatelem
\begin{itemize}
\item např. true negative rate = true negative / healthy
\begin{itemize}
\item = true negative / (true negative + false positive)
\end{itemize}
\item false negative rate = false negative / sick
\begin{itemize}
\item = false negative / (false negative + true positive)
\end{itemize}
\end{itemize}
\item senzitivita \dots{} true positive rate (recall)
\begin{itemize}
\item pravděpodobnost pozitivního výsledku u pozitivního (nemocného) jedince
\end{itemize}
\item specificita \dots{} true negative rate
\begin{itemize}
\item pravděpodobnost negativního výsledku u negativního (zdravého) jedince
\end{itemize}
\end{itemize}
\end{itemize}
\item ohodnocení modelu (testovací data, křížová validace, maximální věrohodnost)
\begin{itemize}
\item testovací data slouží k ověření toho, jak dobře náš model generalizuje
\begin{itemize}
\item snažíme se zjistit, jak dobré výsledky model vrací pro data, která dosud neviděl
\item testovací data nesmím použít v procesu trénování modelu a ladění hyperparametrů
\end{itemize}
\item pokud je potřeba ladit hyperparametry, používá se train-dev-test split
\begin{itemize}
\item parametry modelu trénujeme na trénovacích datech
\item hyperparametry určujeme pomocí vývojových dat (zjistíme, pro které hodnoty hyperparametrů má model nejlepší výkon)
\begin{itemize}
\item vývojová data se někdy označují jako validační (validation set)
\end{itemize}
\item k finálnímu vyhodnocení modelu slouží testovací data
\end{itemize}
\item pokud mám dat málo a nelze z nich jednoduše vyčlenit testovací (nebo vývojová) data, používá se křížová validce (cross-validation)
\begin{itemize}
\item trénovací data rozdělím na $n$ dílů
\item model natrénuju na datech z $n-1$ dílů
\item na zbylých datech model vyhodnotím
\item proces opakuju pro $n$ možných voleb
\item jako výsledné skóre modelu použiju průměr průběžných hodnocení
\end{itemize}
\item maximální věrohodnost (maximum likelihood)
\begin{itemize}
\item máme-li fixní trénovací data, pak likelihood $L(w)$ (kde $w$ jsou váhy modelu) se rovná součinu pravděpodobností targetů trénovacích dat
\item odhadujeme váhy modelu tak, aby byl likelihood (věrohodnost) trénovacích dat co největší
\item přesněji
\begin{itemize}
\item pro model s fixními vahami $w$ máme pravděpodobnostní distribuci $p_\text{model}(x;w)$, že model vygeneruje $x$
\item pro model s fixními vahami $w$ a konkrétní řádek $x$ máme pravděpodobnostní distribuci $p_\text{model}(t\mid x;w)$, jak model klasifikuje $x$
\item místo toho zafixujeme trénovací data a budeme měnit váhy $\to$ dostaneme $L(w)=p_\text{model}(X;w)=\prod_i p_\text{model}(x_i;w)$
\item $w_\text{MLE}=\text{arg max}_w\,p_\text{model}(X;w)$
\begin{itemize}
\item MLE \dots{} maximum likelihood estimation
\end{itemize}
\end{itemize}
\end{itemize}
\end{itemize}
\item přeučení a regularizace (generalizační chyba, včasné zastavení trénování, L2 a L1 regularizace)
\begin{itemize}
\item underfitting - model je příliš slabý (příliš jednoduchý), plně jsme nevyužili jeho potenciál
\item overfitting (přeučení) - model špatně generalizuje, příliš se přizpůsobil trénovacím datům
\item generalizační chyba
\begin{itemize}
\item zachycuje, jak dobře model generalizuje
\item typicky odpovídá výkonu modelu na testovacích datech (případně lze k jejímu určení použít křížovou validaci)
\begin{itemize}
\item výkon = vhodná metrika (F-measure, accuracy, \dots{})
\end{itemize}
\end{itemize}
\item generalizační chybu lze zachytit v grafu vedle trénovací chyby (výkonu modelu na trénovacích datech)
\begin{itemize}
\item pokud se testovací (generalizační) křivka přibližovala k trénovací a v určitý moment se začala vzdalovat, došlo k~přetrénování (overfitting) modelu
\begin{itemize}
\item je potřeba zesílit regularizaci nebo včas zastavit trénování
\end{itemize}
\item pokud se testovací křivka ani nezačala přibližovat, je potřeba trénovat déle
\end{itemize}
\item regularizace
\begin{itemize}
\item intuice
\begin{itemize}
\item penalizujeme modely s většími váhami
\item preferujeme menší změny modelu
\item omezujeme efektivní kapacitu modelu
\item čím je model složitější, tím větší má váhy
\end{itemize}
\item ke ztrátové funkci (loss \textasciitilde{}\textasciitilde{}:.|:;\textasciitilde{}\textasciitilde{}) přičteme normu vah vynásobenou parametrem $\lambda$
\begin{itemize}
\item $L^1$ regularizace \dots{} $\lambda(|w_1|+|w_2|+\dots+|w_n|)$
\item $L^2$ regularizace \dots{} $\lambda(w_1^2+w_2^2+\dots+w_n^2)$
\begin{itemize}
\item nebo také $\lambda\lVert w\rVert^2$, případně $\frac\lambda2\lVert w\rVert^2$, aby nám vyšla hezká derivace
\end{itemize}
\end{itemize}
\item pozor, při regularizaci obvykle vynecháváme váhu, která odpovídá biasu
\end{itemize}
\item jak bojovat s overfittingem
\begin{itemize}
\item $L^2$ regularizací
\item early stoppingem (s trénováním přestaneme, když začne validační chyba růst)
\item vhodným nastavením learning ratu $\alpha$ (případně volbou rozumného learning rate schedule)
\end{itemize}
\end{itemize}
\item prokletí dimenzionality
\begin{itemize}
\item s rostoucím počtem dimenzí (features) potřebujeme čím dál víc (exponenciálně) trénovacích dat, aby model rozumně generalizoval
\item chceme mít v trénovacích datech několik vzorků pro každou kombinaci features $\to$ s~každou další dimenzí potřebujeme násobně víc vzorků
\item k redukci počtu dimenzí se obvykle používají následující přístupy
\begin{itemize}
\item filter
\begin{itemize}
\item nezávislý na modelu, features vybírá na základě jejich vlastností
\item k hodnocení relevance features se používají statistické metody
\begin{itemize}
\item konzistence - nemůžeme odstranit features, pokud bychom tak ztratili možnost odlišit prvky s různými targety (třídami)
\item korelace - dobré features korelují s třídami a nekorelují s jinými features
\item vzdálenost mezi třídami
\item metriky z teorie informace - např. vzájemná informace
\end{itemize}
\end{itemize}
\item wrapper
\begin{itemize}
\item závislý na modelu, features vybírá tak, aby optimalizoval jeho výsledky
\item model se natrénuje, pak se vyhodnotí pomocí standardních technik (accuracy, F-measure apod.)
\end{itemize}
\end{itemize}
\end{itemize}
\end{itemize}
\subsection{Učení založené na příkladech}
\begin{itemize}
\item algoritmus k-nejbližších sousedů
\begin{itemize}
\item regrese: $t=\sum_i\frac{w_i}{\sum_j w_j}\cdot t_i$
\begin{itemize}
\item vážený průměr
\end{itemize}
\item klasifikace
\begin{itemize}
\item pro uniformní váhy můžeme hlasovat (nejčastější třída vyhrává, remízy rozhodujeme náhodně)
\item jinak uvažujeme one-hot kódování $t_i$ opět můžeme použít predikci $t=\sum_i\frac{w_i}{\sum_j w_j}\cdot t_i$ (zvolíme třídu s největší predikovanou pravděpodobností)
\end{itemize}
\item volba vah (weighting)
\begin{itemize}
\item uniform: všech $k$ nejbližších sousedů má stejné váhy
\item inverse: váhy jsou nepřímo úměrné vzdálenosti
\item softmax: váhy jsou přímo úměrné softmaxu záporných vzdáleností
\end{itemize}
\item jak určit nejbližší sousedy / vzdálenosti?
\begin{itemize}
\item $p$-norma: $\|x\|_p=\sqrt[p]{\sum_i|x_i|^p}$
\item obvykle se používá $p\in\set{1,2,3,\infty}$
\item vzdálenost $x$ a $y$ pak lze určit jako $\|x-y\|_p$
\end{itemize}
\end{itemize}
\end{itemize}
\subsection{Lineární regrese}
\begin{itemize}
\item analytické řešení metodou nejmenších čtverců
\begin{itemize}
\item sum of squares error: $\frac12\sum_{i}(x_i^Tw-t_i)^2$
\begin{itemize}
\item lze ho zapsat jako $\frac12\lVert Xw-t\rVert^2$
\end{itemize}
\item snažíme se ho minimalizovat $\to$ hledáme hodnoty, kde je derivace chybové funkce podle každé z vah nulová
\begin{itemize}
\item $\frac{\partial}{\partial w_j}\frac12\sum_{i}(x_i^Tw-t_i)^2=\sum_i x_{ij}(x_i^Tw-t_i)$
\item chceme, aby $\forall j:\sum_i x_{ij}(x_i^Tw-t_i)=0$
\begin{itemize}
\item $X^T(Xw-t)=0$
\item $X^TXw=X^Tt$
\end{itemize}
\end{itemize}
\item je-li $X^TX$ invertibilní, pak lze určit explicitní řešení jako $w=(X^TX)^{-1}X^Tt$
\begin{itemize}
\item přičemž bias reprezentujeme jako součást vah, takže matici $X$ musíme rozšířit o sloupec jedniček
\end{itemize}
\end{itemize}
\item trénování pomocí stochastic gradient descent
\begin{itemize}
\item gradient descent
\begin{itemize}
\item snažíme se minimalizovat hodnotu chybové funkce $E(w)$ pro dané váhy $w$ volbou lepších vah
\begin{itemize}
\item spočítáme $\nabla_wE(w)$ $\to$ tak zjistíme, jak váhy upravit
\item nastavíme $w\leftarrow w-\alpha\nabla_wE(w)$
\begin{itemize}
\item $\alpha$ \dots{} learning rate (hyperparametr), \quotedblbase{}délka kroku\textquotedblleft{}
\end{itemize}
\item tomu se říká gradient descent
\end{itemize}
\item standard gradient descent - používáme všechna trénovací data k výpočtu $\nabla_wE(w)$
\item stochastic (online) gradient descent - používáme jeden náhodný řádek
\item minibatch stochastic gradient descent - používáme $B$ náhodných nezávislých řádků
\end{itemize}
\item algoritmus pro trénink modelu lineární regrese (minibatch SGD)
\begin{itemize}
\item náhodně inicializujeme $w$ (nebo nastavíme na nulu)
\item opakujeme, dokud to nezkonverguje nebo nám nedojde trpělivost
\begin{itemize}
\item samplujeme minibatch řádků s indexy $\mathbb B$
\begin{itemize}
\item buď uniformně náhodně, nebo budeme chtít postupně zpracovat všechna trénovací data (to se obvykle dělá, jednomu takovému průchodu se říká epocha), tedy je nasekáme na náhodné minibatche a procházíme postupně
\end{itemize}
\item $w\leftarrow w-\alpha\frac1{|\mathbb B|}\sum_{i\in \mathbb B}((x_i^Tw-t_i)x_i)-\alpha\lambda w$
\begin{itemize}
\item jako $E$ se používá polovina MSE (s $L^2$ regularizací): $E(w)=\frac1{|\mathbb B|}\sum_{i\in\mathbb B}\frac12 (x_i^Tw-t_i)^2+\frac{\lambda}2 \lVert w\rVert^2$ (stačí to zderivovat)
\end{itemize}
\end{itemize}
\end{itemize}
\end{itemize}
\end{itemize}
\subsection{Logistická regrese}
\begin{itemize}
\item binární klasifikace (sigmoid, metody trénování)
\begin{itemize}
\item aktivační funkce sigmoid \dots{} $\sigma(x)=\frac1{1+e^{-x}}$
\item predikce $y(x;w)=\sigma(\bar y(x;w))=\sigma(x^Tw)$
\begin{itemize}
\item správnou třídu získáme zaokrouhlením
\item přesná hodnota se rovná pravděpodobnosti, že je $x$ klasifikováno jako 1 (uvažujeme třídy 0 a 1)
\item $\bar y$ \dots{} \quotedblbase{}lineární složka\textquotedblleft{} logistické regrese
\end{itemize}
\item velikost vektoru vah $w$ odpovídá počtu features
\begin{itemize}
\item nesmíme zapomenout také na bias, ten opět může být reprezentován jako další váha
\end{itemize}
\item při trénování se jako ztrátová funkce používá negative log likelihood (NLL)
\begin{itemize}
\item pravděpodobnost targetu $t$ pro položku $x$ lze vyjádřit jako $p(t\mid x;w)=\sigma(x^Tw)^{t}(1-\sigma(x^Tw))^{1-{t}}$
\end{itemize}
\item algoritmus trénování pomocí minibatch SGD
\begin{itemize}
\item klasifikujeme do tříd $\set{0,1}$, na vstupu máme dataset $X$ a learning rate $\alpha\in\mathbb R^+$
\item náhodně inicializujeme $w$ (nebo nastavíme na nulu)
\item opakujeme, dokud to nezkonverguje nebo nám nedojde trpělivost
\begin{itemize}
\item samplujeme minibatch řádků s indexy $\mathbb B$
\item $w\leftarrow w-\alpha\frac1{|\mathbb B|}\sum_{i\in \mathbb B}(\sigma(x^Tw)-t)x-\alpha\lambda w$
\begin{itemize}
\item jako $E$ se používá NLL (s $L^2$ regularizací)
\item výsledek se překvapivě podobá lineární regresi
\begin{itemize}
\item gradient je také $(y(x)-t)x$
\end{itemize}
\end{itemize}
\end{itemize}
\end{itemize}
\end{itemize}
\item klasifikace do více tříd (softmax, metody trénování)
\begin{itemize}
\item klasifikujeme do jedné z $K$ tříd
\item parametry modelu: váhová matice $W$ (sloupce odpovídají třídám, řádky featurám), vektor biasů (jeden bias za každou třídu)
\item aktivační funkce \dots{} $\text{softmax}(z)_i=\frac{e^{z_i}}{\sum_je^{z_j}}$
\begin{itemize}
\item softmax je invariantní k přičtení konstanty $\text{softmax}(z+c)_i=\text{softmax}(z)_i$
\item $\sigma(x)=\text{softmax}([x,0])_0=\frac{e^x}{e^x+e^0}=\frac1{1+e^{-x}}$
\end{itemize}
\item klasifikace: $p(C_i\mid x;W)=y(x;W)_i=\text{softmax}(\bar y(x;W))_i=\text{softmax}(x^TW)_i$
\begin{itemize}
\item softmax zajišťuje normalizaci (aby se pravděpodobnosti sečetly na jedničku)
\end{itemize}
\item z linearity $\bar y$ lze dokázat, že oblasti tříd jsou konvexní - nemůže se stát, že by na úsečce mezi dvěma body v jedné třídě byl bod v jiné třídě
\item algoritmus trénování pomocí minibatch SGD
\begin{itemize}
\item klasifikujeme do tříd $\set{0,\dots,K-1}$, na vstupu máme dataset $X$ a learning rate $\alpha\in\mathbb R^+$
\item náhodně inicializujeme $w$ (nebo nastavíme na nulu)
\item opakujeme, dokud to nezkonverguje nebo nám nedojde trpělivost
\begin{itemize}
\item samplujeme minibatch řádků s indexy $\mathbb B$
\item $W\leftarrow W-\alpha\frac1{|\mathbb B|}\sum_{i\in \mathbb B}\left((\text{softmax}(x_i^TW)-1_{t_i})x_i^T\right)^T-\alpha\lambda W$
\begin{itemize}
\item jako $E$ se používá NLL (s $L^2$ regularizací)
\item srovnání gradientu logistické regrese
\begin{itemize}
\item binary: $(y(x)-t)x$
\item multiclass: $((y(x)-1_t)x^T)^T$
\end{itemize}
\item přičemž $1_t$ je one-hot reprezentace targetu $t$ (tedy vektor s jedničkou na $t$-té pozici a nulami všude jinde)
\end{itemize}
\end{itemize}
\end{itemize}
\end{itemize}
\end{itemize}
\subsection{Rozhodovací stromy}
\begin{itemize}
\item algoritmus učení a kritéria větvení
\begin{itemize}
\item v každém vnitřním vrcholu je uložena feature, podle které dělíme, a konkrétní hodnota, která určuje rozhraní
\item každému vrcholu náleží podmnožina trénovacích dat
\begin{itemize}
\item predikce listu odpovídá tomu, jaká data v~něm jsou (při regresi průměrujeme, při klasifikaci volíme největší třídu = většinové hlasování / majority vote)
\end{itemize}
\item možná kritéria větvení ve vrcholu $\mathcal T$
\begin{itemize}
\item squared error criterion (používá se při regresi)
\begin{itemize}
\item squared error criterion říká, jak homogenní je podmnožina trénovacích dat, která náleží vrcholu $\mathcal T$
\item $c_\text{SE}(\mathcal T)=\sum_{i\in I_\mathcal T}(t_i-\hat t_{\mathcal T})^2$
\item $I_\mathcal T$ jsou indexy řádků trénovacích dat, které náleží danému vrcholu
\item $\hat t_\mathcal T$ je průměr targetů v daném vrcholu
\end{itemize}
\item Gini index / Gini impurity (používá se při klasifikaci)
\begin{itemize}
\item $c_\text{Gini}(\mathcal T)=|I_\mathcal T|\sum_kp_\mathcal T(k)(1-p_\mathcal T(k))$
\item jak často by byl náhodně vybraný prvek špatně klasifikován, kdyby byl klasifikován náhodně podle rozdělení $p_\mathcal T$
\begin{itemize}
\item $p_\mathcal T(k)$ \dots{} pravděpodobnost, že vybereme prvek ze třídy $k$
\item $(1-p_\mathcal T(k))$ \dots{} pravděpodobnost, že mu přiřadíme jinou třídu než $k$
\end{itemize}
\item v binární klasifikaci Gini odpovídá squared error ztrátové funkci (skoro MSE)
\end{itemize}
\item entropy criterion (používá se při klasifikaci)
\begin{itemize}
\item $c_\text{entropy}(\mathcal T)=|I_\mathcal T|\cdot H(p_\mathcal T)=-|I_\mathcal T|\sum_k p_\mathcal T(k)\log p_\mathcal T(k)$
\item ze součtu vynecháme třídy s nulovou pravděpodobností, aby nám nerozbily logaritmus
\item entropy criterion lze odvodit ze ztrátové funkce negative log likelihood (NLL)
\end{itemize}
\end{itemize}
\item algoritmus CART (Classification and Regression Trees)
\begin{itemize}
\item trénovací data jsou uložena v listech 
\item začneme s jedním listem
\item list můžeme rozdělit (stane se z něj vnitřní vrchol a pod něj připojíme dva nové listy)
\item list dělíme tak, aby rozdíl $c_{\mathcal T_L}+c_{\mathcal T_R}-c_{\mathcal T}$ byl co nejmenší
\begin{itemize}
\item $c_\mathcal T$ \dots{} criterion současného vrcholu
\item $c_{\mathcal T_L}$ \dots{} criterion nového levého syna
\item $c_{\mathcal T_R}$ \dots{} criterion nového pravého syna
\end{itemize}
\item někdy stanovujeme omezení: maximální hloubku stromu, minimální počet dat v~děleném listu (abychom nedělili zbytečně listy, co mají málo dat), maximální počet listů
\begin{itemize}
\item obvykle se listy dělí prohledáváním do hloubky
\item pokud je stanoven maximální počet listů, je vhodnější je dělit ty s nejmenším rozdílem $c_{\mathcal T_L}+c_{\mathcal T_R}-c_{\mathcal T}$
\end{itemize}
\end{itemize}
\item dosud nás vždycky zajímaly parametry, které minimalizují ztrátovou funkci
\item tím, že minimalizujeme dělicí kritérium, minimalizujeme minimální dosažitelnou hodnotu ztrátové funkce pro daný strom
\end{itemize}
\item prořezávání
\begin{itemize}
\item prořezávání (pruning) lze použít ke zjednodušení rozhodovacího stromu, tím je možné omezit overfitting
\item prořezávat lze shora dolů nebo zespoda nahoru (bottom-up pruning $\times$ top-down pruning)
\item příklad: reduced error pruning
\begin{itemize}
\item bottom-up přístup
\item na testovacích datech ověřujeme přesnost (accuracy)
\item postupně zkoušíme rušit větvení ve vnitřních vrcholech - pokud tím nedojde ke snížení přesnosti, tak změnu provedeme
\begin{itemize}
\item tedy se z vnitřního vrcholu stane list, který predikuje třídu s největším počtem položek
\end{itemize}
\end{itemize}
\end{itemize}
\end{itemize}
\subsection{Metoda podpůrných vektorů}
\begin{itemize}
\item klasifikátor pro lineárně separabilní třídy
\begin{itemize}
\item hard-margin SVM (support-vector machine)
\item je to kernelová (jádrová) metoda strojového učení
\begin{itemize}
\item model má k dispozici spoustu polynomiálních features, aniž bychom je museli explicitně počítat
\item používáme jádrovou funkci $K$
\end{itemize}
\item model je neparametrický - využívá se duální formulace
\begin{itemize}
\item nemáme váhy, ale koeficienty u trénovacích dat
\item duální formulace vychází z pozorování, že klasické váhy jsou lineárními kombinacemi trénovacích dat
\end{itemize}
\item princip je podobný jako u perceptronu
\begin{itemize}
\item jedná se o binární klasifikaci do tříd $\set{-1,1}$
\item tentokrát ale chceme zajistit, aby byla dělicí nadrovina (decision boundary) co nejlepší - aby byl margin (mezera mezi nadrovinou a body) co největší
\end{itemize}
\item jak se odvozuje
\begin{itemize}
\item uvažujeme váhy $w$ a bias $b$, dále mapování $\varphi$ do prostoru features, predikci označíme $y(x)=\varphi(x)^Tw+b$
\item pro všechny body chceme maximalizovat vzdálenost od decision boundary
\begin{itemize}
\item váhy i bias můžeme libovolně škálovat, proto určíme, že pro body nejbližší k decision boundary bude platit $t_iy(x_i)=1$
\item díky tomu nám stačí \quotedblbase{}minimalizovat váhy\textquotedblleft{} takto: $\mathrm{argmin}_{w,b}\,\frac12\|w\|^2$
\begin{itemize}
\item za předpokladu, že $t_iy(x_i)\geq 1$
\end{itemize}
\end{itemize}
\item tuhle minimalizaci s omezující podmínkou provádíme pomocí lagrangiánu a KKT podmínek
\begin{itemize}
\item místo multiplikátorů $\lambda_i$ se tradičně používají $a_i$
\end{itemize}
\item díky KKT podmínkám víme, že nás při maximalizaci lagrangiánu zajímají jenom některá trénovací data - ta, která leží na okraji (support vektory)
\begin{itemize}
\item je tam podmínka $a_i(t_iy(x_i)-1)=0$
\item z toho vyplývá, že pro každé $x_i$ platí $t_iy(x_i)=1$ (bod je na marginu) nebo $a_i=0$ (bod se nepoužívá k predikci)
\end{itemize}
\end{itemize}
\item predikce \dots{} $y(x)=\sum_i a_it_iK(x,x_i)+b$
\end{itemize}
\item klasifikátor pro lineárně neseparabilní třídy
\begin{itemize}
\item soft-margin SVM
\item v mnohém se podobá hard-margin SVM
\item dovolíme některým bodům, aby porušily pravidlo (aby byly uvnitř marginu nebo na opačné straně decision boundary)
\item pořídíme si slack variables $\xi_i\geq 0$
\begin{itemize}
\item pro body, které porušují $t_iy(x_i)\geq1$ bude platit $\xi_i=|t_i-y(x_i)|$
\item jinak $\xi_i=0$
\end{itemize}
\item takže podmínku $t_iy(x_i)\geq1$ nahradíme $t_iy(x_i)\geq1-\xi_i$
\item úpravy lagrangiánu
\begin{itemize}
\item musíme tam přidat další multiplikátory, které zajistí nezápornost $\xi_i$
\item výsledný lagrangián bude stejný jako minule, akorát $a_i$ (odpovídají lambdám) budou navíc omezeny shora nějakým $C$
\end{itemize}
\item support vektory teďka nebudou jenom body na okrajích marginu, ale i všechny, které mají kladné $\xi_i$ (jsou divné)
\item pozorování: $\xi_i=\max(0,1-t_iy(x_i))$
\begin{itemize}
\item téhle funkci se říká hinge loss
\end{itemize}
\item soft-margin SVM lze vlastně formulovat jako $\text{argmin}_{w,b}\,C\sum_i\mathcal L_\mathrm{hinge}(t_i,y(x_i))+\frac12\|w\|^2$
\begin{itemize}
\item to nám hodně připomíná klasický vzorec pro logistickou regresi apod.
\item akorát místo regularizační konstanty $\lambda$ tady máme $C$, které hraje opačnou roli (čím větší $C$, tím slabší regularizace)
\end{itemize}
\item k trénování se používá sequential minimal optimization
\begin{itemize}
\item bereme řádky po jednom, k tomu $a_i$, náhodně $a_j$ a taky bias, postupně zlepšujeme lagrangián
\end{itemize}
\end{itemize}
\item jádrové funkce
\begin{itemize}
\item idea
\begin{itemize}
\item uvažujeme duální formulaci modelu (místo vah máme koeficienty, kterými přispívají řádky trénovacích dat)
\item máme polynomiální features (pomocí zobrazení $\varphi$)
\item při predikci bychom museli $N$-krát počítat $\varphi(x_i)^T\varphi(z)$
\item když si ale například představíme features 0. až 3. řádu, tak si můžeme všimnout toho, že $\varphi(x)^T\varphi(z)=1+x^Tz+(x^Tz)^2+(x^Tz)^3$
\end{itemize}
\item zobecnění
\begin{itemize}
\item budeme uvažovat kernelové funkce (kernely), které mají všechny následující vzorec (pro různé $\varphi$): $K(x,z)=\varphi(x)^T\varphi(z)$
\item kernel je tedy funkce, která dostane dva vektory a ona mi vrátí součin features těch vektorů
\item výše jsme si jeden takový kernel ukázali
\end{itemize}
\item polynomiální kernel stupně $d$ (homogenní polynomiální kernel)
\begin{itemize}
\item vyrábí kombinace $d$ vstupních features
\item $K(x,z)=(\gamma x^Tz)^d$
\item $\gamma$ je hyperparametr, který škáluje features
\end{itemize}
\item polynomiální kernel stupně nejvýše $d$ (nehomogenní polynomiální kernel)
\begin{itemize}
\item $K(x,z)=(\gamma x^Tz+1)^d$
\end{itemize}
\item Gaussian Radial basis function (RBF) kernel
\begin{itemize}
\item obsahuje polynomiální features všech řádů
\item $K(x,z)=e^{-\gamma\|x-z\|^2}$
\item je to funkce vzdálenosti
\begin{itemize}
\item pro stejné vektory vrací jedničku, pro hodně vzdálené vektory vrací nulu, pro ostatní něco mezi tím
\item dá se to interpretovat jako rozšířenou verzi algoritmu $k$ nejbližších sousedů ($k$-NN)
\item $\gamma$ ovlivňuje, co už je \quotedblbase{}moc daleko\textquotedblleft{}
\end{itemize}
\item polynomiální features vysokých řádů mají nízké váhy
\end{itemize}
\item kernely se typicky používají tam, kde potřebujeme lineárně separabilní data (například v SVM) - transformují nám původní data do více dimenzí, tam už je pak obvykle můžeme rozumně separovat
\end{itemize}
\item klasifikace do více tříd
\begin{itemize}
\item kombinujeme několik binárních klasifikátorů
\item první možný (nepoužívaný!) přístup: one-versus-rest (ovr) schéma
\begin{itemize}
\item natrénuju $K$ nezávislých binárních klasifikátorů (patří prvek konkrétní třídě? ano/ne)
\item při predikci vyberu třídu, které její klasifikátor přisuzuje největší pravděpodobnost
\item je potřeba klasifikátory kalibrovat, aby se ty pravděpodobnosti daly porovnávat!
\item takhle se to u SVM nedělá
\end{itemize}
\item alternativní přístup: one-versus-one schéma
\begin{itemize}
\item klasifikátor na každou dvojici tříd, které existují
\item většinové hlasování (každý klasifikátor hlasuje pro jednu třídu)
\item klasifikátorů je víc, ale trénujeme je na méně datech
\end{itemize}
\end{itemize}
\end{itemize}
\subsection{Kombinace více modelů}
\begin{itemize}
\item bagging a boosting
\begin{itemize}
\item jedná se o dva různé přístupy k trénování více modelů
\item bagging
\begin{itemize}
\item modely se trénují nezávisle
\item pro každý model náhodně vybereme vzorek trénovacích dat, na nichž se učí (typicky samplujeme s opakováním = bootstrapping)
\end{itemize}
\item boosting
\begin{itemize}
\item modely se trénují sekvenčně
\item snažíme se postupně opravovat chyby předchozích modelů
\item princip algoritmu AdaBoost
\begin{itemize}
\item postupně trénujeme slabé klasifikátory (např. mělké rozhodovací stromy/pařezy) na vážených trénovacích datech
\item v každé iteraci upravujeme váhy trénovacích dat, aby špatně klasifikovaná data měla větší váhy
\item rovněž přiřazujeme váhy klasifikátorům podle jejich přesnosti (accuracy)
\end{itemize}
\end{itemize}
\end{itemize}
\item metoda náhodných lesů
\begin{itemize}
\item trénování
\begin{itemize}
\item pro každý strom náhodně (s opakováním) vybereme sample trénovacích dat (=~bagging)
\item pro každý vrchol náhodně vybereme podmnožinu features, které při dělení bereme v úvahu
\end{itemize}
\item predikce
\begin{itemize}
\item u regrese průměrujeme hodnoty jednotlivých stromů
\item u klasifikace hlasujeme
\end{itemize}
\item les vs. strom
\begin{itemize}
\item strom se dá rychle natrénovat, jeho predikci lze snadno interpretovat (jako sérii rozhodnutí)
\item les je obvykle přesnější a lépe generalizuje
\end{itemize}
\end{itemize}
\end{itemize}
\subsection{Statistické testy}
\begin{itemize}
\item studentův t-test (jednovýběrový a dvouvýběrový)
\begin{itemize}
\item jednovýběrový t-test
\begin{itemize}
\item způsob, jak získat intervalový odhad pro $\theta$, neznáme-li rozptyl $\sigma^2$
\item uvažujeme statistiku $T=\frac{\overline{X_n}-\theta_0}{\bar\sigma/\sqrt n}$, která má t-rozdělení s $n-1$ stupni volnosti, pokud $H_0$ platí
\item intervalový odhad
\begin{itemize}
\item najdeme bodový odhad $\hat\theta$ pro $\theta$
\begin{itemize}
\item třeba pokud nás zajímá $\mu$, tak prostě použijeme výběrový průměr
\end{itemize}
\item máme $n$ hodnot
\item $\delta=\frac{\bar\sigma}{\sqrt n}\cdot \Psi^{-1}_{n-1}(1-\alpha/2)$
\begin{itemize}
\item kde $\Psi^{-1}_{n-1}$ je kvantilová funkce studentova t-rozdělení
\end{itemize}
\item vrátíme $[\hat\theta-\delta,\hat\theta+\delta]$
\begin{itemize}
\item v tomto intervalu bude pravděpodobně reálná hodnota $\theta$ pro populaci, z~níž jsme data vybrali
\end{itemize}
\end{itemize}
\item testování hypotézy
\begin{itemize}
\item např. pokud ověřujeme, že se střední hodnota rovná nějaké konstantě, tak konstantu i výběrový průměr dosadíme do vzorce pro $T$
\item $H_0$ \dots{} $T$ má t-rozdělení s $n-1$ stupni volnosti
\item $H_0$ zamítáme, když je spočítaná hodnota statistiky větší než kritická hodnota
\end{itemize}
\item poznámka: proč uvažujeme $n-1$ stupňů volnosti, když máme $n$ hodnot?
\begin{itemize}
\item z $n-1$ hodnot můžu $n$-tou hodnotu dopočítat
\end{itemize}
\end{itemize}
\item dvouvýběrový t-test
\begin{itemize}
\item můžeme testovat třeba to, jestli se rovnají střední hodnoty parametru $\theta$ u dvou různých populací (samplujeme nezávisle, vzorků může být různý počet)
\item např. $H_0$ \dots{} $\mu_X=\mu_Y$
\item opět uvažujeme testovou statistiku $T$ s t-rozdělením
\begin{itemize}
\item $T=\frac{\overline{X_n}-\overline{Y_m}}{\mathrm{se}}$
\item přičemž $\mathrm{se}$ je směrodatná chyba
\item $\mathrm{se}=\sqrt{\frac{\sigma_X^2}{n}+\frac{\sigma_Y^2}{m}}$
\end{itemize}
\item případně opět $\delta=\mathrm{se}\cdot \Phi^{-1}(1-\alpha/2)$
\end{itemize}
\item bonus: párový test
\begin{itemize}
\item data jsou ve dvojicích, dvojice dat spolu nějak souvisí
\item např. $H_0$ \dots{} $\mu_X=\mu_Y$
\item uvažuju $R_i=Y_i-X_i$ (rozdíl dvojice), použiju jednovýběrový test
\end{itemize}
\end{itemize}
\item chí-kvadrát test (test dobré shody)
\begin{itemize}
\item Pearsonův chí-kvadrát test dobré shody
\item $O_i$ \dots{} pozorovaná četnost kategorie $i$
\item $E_i$ \dots{} očekávaná četnost kategorie $i$
\item testová statistika: $\sum_i\frac{(E_i-O_i)^2}{E_i}$
\item typická úloha: je kostka spravedlivá?
\begin{itemize}
\item použijeme $\chi^2$ distribuci pro 5 stupňů volnosti
\begin{itemize}
\item pokud $\chi^2$ pro naši konkrétní kostku náleží kritickému oboru $W$, prohlásíme, že kostka není spravedlivá
\end{itemize}
\end{itemize}
\item $H_0$ \dots{} pozorované četnosti jednotlivých kategorií odpovídají jejich očekávaným četnostem
\item pokud $H_0$ platí, má testová statistika rozdělení $\chi^2$ rozdělení s $n-1$ stupni volnosti, přičemž $n$ je počet kategorií
\item $H_0$ zamítáme, když je spočítaná hodnota statistiky větší než kritická hodnota
\end{itemize}
\end{itemize}
\subsection{Učení bez učitele}
\begin{itemize}
\item shlukování (algoritmus k-means)
\begin{itemize}
\item algoritmus
\begin{itemize}
\item na vstupu máme body $x_1,\dots,x_N$, počet clusterů $K$
\item inicializujeme středy clusterů $\mu_1,\dots,\mu_K$
\begin{itemize}
\item buď náhodně (ze vstupních bodů)
\item nebo pomocí \texttt{kmeans++}
\begin{itemize}
\item první střed se vybere náhodně (ze vstupních bodů)
\item každý další střed vybereme z nepoužitých bodů tak, že pravděpodobnost výběru každého bodu úměrně roste s druhou mocninou vzdálenosti jeho nejbližšího středu (clusteru)
\end{itemize}
\end{itemize}
\item opakujeme, dokud to nezkonverguje nebo nám nedojde trpělivost
\begin{itemize}
\item body přiřadíme clusterům (každý bod přiřadíme tomu clusteru, k jehož středu má nejblíž)
\item aktualizujeme středy clusterů (vždy zprůměrujeme body, které jsme danému clusteru přiřadili)
\end{itemize}
\end{itemize}
\item použití
\begin{itemize}
\item $K$-means se používá ke clusteringu, tedy k dělení dat do skupin, clusterů
\item je to technika učení bez učitele (unsupervised), neznáme targety
\item konkrétní příklady: seskupení pixelů podobné barvy, určení skupin zákazníků
\end{itemize}
\item ztrátová funkce, kterou algoritmus optimalizuje
\begin{itemize}
\item $J=\sum_{i=1}^N\sum_{k=1}^Kz_{i,k}\|x_i-\mu_k\|^2$
\item kde $z_{i,k}$ je indikátor přiřazení bodu $x_i$ do clusteru $k$
\end{itemize}
\item konvergence
\begin{itemize}
\item aktualizace přiřazení bodů ke clusterům nezvyšuje loss
\item aktualizace středů clusterů nezvyšuje loss
\item takže $K$-means konverguje k lokálnímu optimu
\end{itemize}
\end{itemize}
\item hierarchické shlukování
\begin{itemize}
\item dva přístupy: postupné spojování menších shluků (aglomerativní shlukování) nebo naopak dělení velkých shluků na menší (divisivní shlukování) podle zadaných pravidel
\item algoritmus spojování menších clusterů/shluků
\begin{itemize}
\item na začátku je každý cluster singleton (patří tam jeden bod)
\item slučujeme nejpodobnější clustery, dokud nezískáme cílový počet clusterů
\end{itemize}
\item dají se používat různé metriky podobnosti shluků (např. vzdálenost mezi centroidy shluků nebo přímo mezi jejich body, průměry vzdáleností, \dots{})
\item příklady algoritmů divisivního shlukování
\begin{itemize}
\item DIANA (divisive analysis)
\begin{itemize}
\item na začátku jsou všechny body v jednom clusteru
\item v něm najdeme bod, který se nejvíc liší od všech ostatních a oddělíme ho do samostatného clusteru
\item do téhož (nového) clusteru po jednom přidáváme i další body z původního clusteru, pokud se více hodí do nového clusteru (jsou mu blíže než původnímu)
\item v každé fázi algoritmu se jeden cluster rozdělí na dva (ty jsou jakoby vnořené do původního clusteru)
\item algoritmus končí, jakmile je každý bod ve vlastním clusteru
\item můžeme získat dendrogram - graf, jak se clustery postupně dělily
\begin{itemize}
\item když ho vhodně zařízneme, tak získáme použitelné dělení (aby clusterů nebylo ani málo, ani moc)
\end{itemize}
\end{itemize}
\item principal directions divisive partitioning (PDDP)
\begin{itemize}
\item používá princip PCA - dělí clustery pomocí projekce na hlavní komponenty
\end{itemize}
\end{itemize}
\item většinou se jedná o hladové algoritmy $\to$ nelze zaručit, že je řešení optimální
\end{itemize}
\item redukce dimenzionality (analýza hlavních komponent)
\begin{itemize}
\item singulární rozklad (SVD, singular value decomposition)
\begin{itemize}
\item mějme matici $X\in\mathbb R^{m\times n}$ s hodností $r$
\item $X=U\Sigma V^T$
\begin{itemize}
\item $U\in\mathbb R^{m\times m}$ ortonormální
\item $\Sigma\in\mathbb R^{m\times n}$ diagonální matice s nezápornými hodnotami (tzv. singulárními čísly) zvolenými tak, aby byly uspořádány sestupně
\item $V\in\mathbb R^{n\times n}$ ortonormální
\end{itemize}
\item možný pohled na dekompozici
\begin{itemize}
\item $X=\sigma_1u_1v_1^T+\sigma_2u_2v_2^T+\dots+\sigma_ru_rv_r^T$
\item $\sigma_i$ \dots{} $i$-té singulární číslo
\item $u_i$ \dots{} $i$-tý sloupec matice $U$
\item $v_i^T$ \dots{} $i$-tý řádek matice $V^T$
\end{itemize}
\item redukovaná verze SVD
\begin{itemize}
\item $\sigma$ označme vektor singulárních čísel
\item pro matici $X$ s hodností $r$ bude v tom vektoru nejvýše $r$ nenulových hodnot
\item tedy matici $\Sigma$ můžeme redukovat na rozměry $r\times r$, podobně $U$ můžeme redukovat na $m\times r$ a $V^T$ na $r\times n$, touto kompresí nepřijdeme o žádné informace
\end{itemize}
\item dokonce můžeme $\Sigma$ zmenšit ještě víc
\begin{itemize}
\item zvolíme $k\lt r$
\item necháme jenom $k$ největších singulárních čísel, všechny ostatní zahodíme
\item tím už se nějaké informace ztratí
\item budeme mít aproximaci původní matice $X$, ta aproximace bude mít hodnost $k$
\end{itemize}
\end{itemize}
\item algoritmus určení PCA (principal component analysis) $M$-té dimenze na základě datové matice $X$
\begin{itemize}
\item spočteme SVD matice $(X-\bar x)$
\begin{itemize}
\item kde $\bar x$ je průměr řádků matice (vektor průměrných hodnot jednotlivých features)
\item k určení SVD (respektive nalezení $V$) lze použít algoritmus power iteration
\begin{itemize}
\item do proměnné $S$ uložíme kovarianční matici $\text{cov}(X)=\frac 1N(X-\bar x)^T(X-\bar x)$
\item pro každou komponentu hledáme vektor $v_i$ (sloupec matice $V$) tak, že začneme s náhodným vektorem a poté ho zleva násobíme maticí $S$ a normalizujeme (dělíme vlastní normou), dokud to nezkonverguje
\item od matice $S$ na konci každého kroku odečteme $\lambda_iv_iv_i^T$, kde $\lambda_i$ je norma vektoru $v_i$ před poslední normalizací
\end{itemize}
\end{itemize}
\item necháme prvních $M$ řádků $V^T$, ostatní zahodíme (takže $V$ bude mít $M$ sloupců)
\item pokud chceme získat hodnoty nových $M$ komponent, tak stačí vynásobit $XV$, tak dostaneme matici $N\times M$
\item funguje to díky tomu, že pomocí SVD matice $(X-\bar x)$ získáváme vlastní vektory kovarianční matice $\text{cov}(X)=\frac 1N(X-\bar x)^T(X-\bar x)$
\end{itemize}
\item aplikace PCA
\begin{itemize}
\item aproximace matic
\item snížení šumu v datech
\item redukce dimenzí dat, extrakce features
\item vizualizace dat
\item praktické použití: doporučovací systémy
\end{itemize}
\item příklad použití PCA
\begin{itemize}
\item \href{https://towardsdatascience.com/pca-clearly-explained-how-when-why-to-use-it-and-feature-importance-a-guide-in-python-7c274582c37e/}{[obr.: příklad]}
\item PCA postupně hledá osy s největším rozptylem
\end{itemize}
\end{itemize}
\end{itemize}
